\documentclass[12pt]{article}
\usepackage{amsmath, amssymb}
\usepackage[margin=1in]{geometry}
\setlength{\parskip}{6pt}
\title{QUBO Formulation: Cardinality-Constrained Portfolio Selection Problem}
\date{}
\begin{document}
\maketitle

\subsection*{Cardinality-Constrained Portfolio Selection Problem}

\paragraph{Problem Statement.}
Given a set of $N$ assets, each asset $i$ is associated with an expected return score $r_i$. Each distinct pair of assets $(i,j)$ is associated with a covariance or risk penalty $c_{ij}$. The task is to select exactly $K$ assets to form a portfolio. The objective is to maximize the risk-adjusted score, defined as the sum of the expected returns of the selected assets minus the sum of the pairwise risk penalties between the selected assets.

\paragraph{Decision Variables.}
For each asset $i \in \{0, 1, \dots, N-1\}$, we introduce a binary decision variable $x_i \in \{0,1\}$. The variable $x_i = 1$ indicates that asset $i$ is included in the portfolio, while $x_i = 0$ indicates that it is excluded.

\paragraph{QUBO Derivation.}
\textbf{Step 1.} The original problem seeks to maximize the risk-adjusted score. To conform to the standard QUBO minimization convention, we negate the objective function:
\begin{align}
  \min_{x} \left( -\sum_{i=0}^{N-1} r_i x_i + \sum_{0 \le i < j \le N-1} c_{ij} x_i x_j \right).
\end{align}

\textbf{Step 2.} The cardinality constraint requires that exactly $K$ assets are selected:
\begin{align}
  \sum_{i=0}^{N-1} x_i = K.
\end{align}

\textbf{Step 3.} We enforce the constraint via a quadratic penalty term with coefficient $P > 0$:
\begin{align}
  \mathcal{P}(x) = P \left( \sum_{i=0}^{N-1} x_i - K \right)^2.
\end{align}
Expanding the square and applying the idempotent property $x_i^2 = x_i$ for binary variables yields:
\begin{align}
  \left( \sum_{i=0}^{N-1} x_i - K \right)^2 &= \sum_{i=0}^{N-1} x_i^2 + 2 \sum_{0 \le i < j \le N-1} x_i x_j - 2K \sum_{i=0}^{N-1} x_i + K^2 \\
  &= \sum_{i=0}^{N-1} (1 - 2K) x_i + 2 \sum_{0 \le i < j \le N-1} x_i x_j + K^2.
\end{align}
Multiplying by $P$ gives the general penalty contribution:
\begin{align}
  \mathcal{P}(x) = P \sum_{i=0}^{N-1} (1 - 2K) x_i + 2P \sum_{0 \le i < j \le N-1} x_i x_j + P K^2.
\end{align}

\textbf{Step 4.} Combining the negated objective and the penalty term yields the total QUBO Hamiltonian $H(x)$:
\begin{align}
  H(x) = \sum_{i=0}^{N-1} \left( P(1-2K) - r_i \right) x_i + \sum_{0 \le i < j \le N-1} \left( c_{ij} + 2P \right) x_i x_j + P K^2.
\end{align}

\textbf{Step 5.} Reading off the coefficients from $H(x)$, the upper-triangular QUBO matrix entries $Q_{i,j}$ and the constant offset $c$ are given by:
\begin{align}
  Q_{i,i} &= P(1-2K) - r_i, & \text{for } i \in \{0, \dots, N-1\}, \\
  Q_{i,j} &= c_{ij} + 2P, & \text{for } 0 \le i < j \le N-1, \\
  Q_{i,j} &= 0, & \text{for } i > j, \\
  c &= P K^2.
\end{align}

\paragraph{Penalty Weight Justification.}
The penalty coefficient $P$ must be chosen sufficiently large to ensure that any violation of the cardinality constraint incurs a higher energy cost than the maximum possible improvement in the objective function. A safe lower bound is $P > \max_{i} |r_i| + \sum_{0 \le i < j \le N-1} |c_{ij}|$. This guarantees that the penalty term strictly dominates the objective term for any infeasible assignment, thereby forcing the global minimum to reside within the feasible region.

\paragraph{Correctness Argument.}
Minimizing $H(x)$ over $\{0,1\}^N$ recovers the optimal portfolio because (i) all feasible assignments satisfying $\sum_{i} x_i = K$ yield a penalty term of exactly zero, leaving only the original risk-adjusted objective to be minimized; and (ii) any infeasible assignment incurs a penalty strictly greater than the maximum possible gain achievable by altering the objective function. Consequently, the global minimum of $H(x)$ must correspond to a feasible assignment that optimally balances expected returns and risk penalties.

\paragraph{Illustrative Example.}
Consider a concrete instance with $N=3$ assets and a target cardinality $K=2$. The parameters are $r_0, r_1, r_2$ for expected returns and $c_{01}, c_{02}, c_{12}$ for pairwise risk penalties. Substituting $K=2$ into the general formulas, the concrete objective becomes:
\begin{align}
  \min_{x} \left( -r_0 x_0 - r_1 x_1 - r_2 x_2 + c_{01} x_0 x_1 + c_{02} x_0 x_2 + c_{12} x_1 x_2 \right).
\end{align}
The penalty expansion yields linear coefficients $P(1-4) = -3P$, quadratic coefficients $2P$, and an offset $4P$. The resulting Hamiltonian is:
\begin{align}
  H(x) &= (-3P - r_0)x_0 + (-3P - r_1)x_1 + (-3P - r_2)x_2 \nonumber \\
  &\quad + (c_{01} + 2P)x_0 x_1 + (c_{02} + 2P)x_0 x_2 + (c_{12} + 2P)x_1 x_2 + 4P.
\end{align}
The corresponding QUBO matrix entries and offset are:
\begin{align}
  Q_{0,0} &= -3P - r_0, & Q_{1,1} &= -3P - r_1, & Q_{2,2} &= -3P - r_2, \\
  Q_{0,1} &= c_{01} + 2P, & Q_{0,2} &= c_{02} + 2P, & Q_{1,2} &= c_{12} + 2P, \\
  c &= 4P.
\end{align}
Since $K=2$, there are exactly three feasible assignments: $x^{(1)} = (1,1,0)$, $x^{(2)} = (1,0,1)$, and $x^{(3)} = (0,1,1)$. Evaluating $H(x)$ at these points (where the penalty term vanishes) gives:
\begin{align}
  H(x^{(1)}) &= -r_0 - r_1 + c_{01} + 4P, \\
  H(x^{(2)}) &= -r_0 - r_2 + c_{02} + 4P, \\
  H(x^{(3)}) &= -r_1 - r_2 + c_{12} + 4P.
\end{align}
The optimal portfolio corresponds to the assignment $x^*$ that minimizes $H(x^*)$. For instance, if $H(x^{(1)})$ is the minimum among the three, the optimal solution is $x^* = (1,1,0)$ with energy $H(x^*) = -r_0 - r_1 + c_{01} + 4P$. The maximum risk-adjusted score for this instance is recovered as $-(H(x^*) - c) = r_0 + r_1 - c_{01}$.
\end{document}